\chapter{Data dictionary}\label{app:dict}

\begin{longtable}{@{}p{3.1cm}p{7.4cm}p{3.6cm}@{}}
\caption{Analysis variables of the mental-health survey (after cleaning).}\label{tab:dict}\\
\toprule
Variable & Definition (source item) & Type / levels \\
\midrule\endfirsthead
\toprule Variable & Definition (source item) & Type / levels \\ \midrule\endhead
\bottomrule\endfoot
\texttt{dep} & Probable depression: PHQ-9 total $\geq$ 10 (Q27\_1--9) & Binary (outcome) \\
\texttt{anx} & Probable anxiety: BAI total $\geq$ 16 (Q24\_1--21) & Binary (outcome) \\
\texttt{depression} & PHQ-9 total & 0--27 \\
\texttt{anxiety} & BAI total & 0--63 \\
\texttt{stress} & PSS-4 total, authors' scoring (Q26\_1--4) & 0--16 \\
\texttt{optimism} & LOT-R optimism items 1, 4, 10 (Q25) & 0--12 \\
\texttt{pessimism} & LOT-R pessimism items 3, 7, 9 (Q25) & 0--12 \\
\texttt{mindfulness} & MAAS-5 total; higher = more mindful (Q29\_1--5) & 5--30 \\
\texttt{sex} & Participant's sex (Q2); third gender ($n=2$) excluded & Male, Female \\
\texttt{age} & Age in years (Q3) & Continuous \\
\texttt{education} & Highest education (Q4), collapsed & 4 levels \\
\texttt{mh\_history} & History of mental disturbance (Q5) & No, Yes, Prefer not to say \\
\texttt{division} & Administrative division (Q6) & 8 levels \\
\texttt{residence} & Rural or urban (Q7) & 2 levels \\
\texttt{living} & Living arrangement (Q8), collapsed & 4 levels \\
\texttt{profession} & Profession (Q9) & 4 levels \\
\texttt{frontliner} & Serving as a frontline worker (Q10) & Binary \\
\texttt{ses} & Subjective socio-economic ladder (Q11) & 0--10 \\
\texttt{married} & Marital status (Q13) & Binary \\
\texttt{children} & Has children (Q14; unmarried coded No) & Binary \\
\texttt{covid\_self} & Own infection (Q17) and treatment place (Q19) & Not infected, home, hospital \\
\texttt{covid\_family} & Family infection (Q20) and death (Q21) & None, infected no death, a death \\
\texttt{vaccine} & Vaccine doses received (Q22) & None, one, two \\
\texttt{phq2} & Sum of PHQ-9 items 1 and 2 (Q27\_1--2); reference ceiling in the ablation only & 0--6 (not a predictor) \\
\texttt{suspect} & Data-quality flag: responses to the 19 non-outcome items (all ten LOT-R items including fillers, PSS-4, MAAS-5) identical to those of $\geq$ 4 other respondents & Binary (not a predictor) \\
\end{longtable}

\chapter{Supplementary tables}\label{app:supp}
\begingroup\setlength{\LTcapwidth}{\linewidth}
\tabin{tab_table1}
\tabin{tab_or}
\endgroup
\clearpage
\tabin{tab_panel_attrition}
\tabin{tab_ml_tuning}
\tabin{tab_ml_anxiety}
\tabin{tab_ml_crosscheck}

\chapter{Reproducibility}\label{app:repro}

\section{Project structure}
Everything under \texttt{output/} and \texttt{data/processed/} is regenerated by \texttt{reproduce.sh} and is not stored in the repository. The raw files in \texttt{data/raw/} can be re-created with \texttt{scripts/download\_data.sh}.
\begin{lstlisting}
analysis-project/
|-- reproduce.sh                 # one command: download -> analysis -> PDF
|-- covid_bd.Rproj
|-- data/
|   |-- raw/                   # downloads (OWID, OxCGRT, Google, Mendeley)
|   `-- processed/             # cleaned data and model objects (generated)
|-- scripts/
|   |-- download_data.sh       # fetches raw data, keeps Bangladesh rows
|   |-- 00_setup.R             # packages, theme, table and macro helpers
|   |-- 01_national.R          # national: waves, PELT, PCA, differenced model
|   |-- 02_panel.R             # panel: attrition, paired tests, quintiles, ANCOVA
|   |-- 03_mh_explore.R        # survey: cleaning, quality flag, aOR/aPR
|   |-- 04_mh_profiles.R       # survey: Gaussian mixtures and their checks
|   |-- 05_mh_ml.R             # survey: fixed and repeated splits, ablation
|   `-- 06_ml_crosscheck.R     # R vs Python comparison table
|-- python/ml_crosscheck.py    # scikit-learn refit of the tier-B models
|-- output/                    # figures/, tables/, macros/ (generated)
`-- report/                    # main.tex, chapters/, references.bib
\end{lstlisting}

\section{Feature whitelist}
Excerpt from \texttt{scripts/05\_mh\_ml.R}:
\begin{lstlisting}[language=R]
TIER <- list(
  A = c("sex", "age", "education", "mh_history", "division", "residence",
        "living", "profession", "frontliner", "ses", "married", "children", "covid_self",
        "covid_family", "vaccine"),
  B = c("optimism", "pessimism", "mindfulness"),
  C = c("anxiety", "stress"))
TIER$B <- c(TIER$A, TIER$B); TIER$C <- c(TIER$B, TIER$C)
stopifnot(!any(c("depression", "dep", "phq2") %in% unlist(TIER)))
\end{lstlisting}
